\chapter{Co-expression with directional extensions} \label{ch:functions}

When directional meaning is defined as path information attached to a motion event expressed elsewhere in an utterance, the result is a semantically constrained distribution of directional meaning. However, in Chadic languages, as elsewhere, morphemes that express directional meaning are not morphosyntactically restricted to only occurring in an utterance that describes a motion event. When directional extensions are used in contexts where the predicate is not describing a motion event, there are a number of non-directional functions that the verbal extension may express. This chapter examines the non-directional functions that are co-expressed with directional meaning in verbal extensions, adding to a growing body of literature on this topic covering languages from around the world \citep{bourdinMarkingDirectionalDeixis2006,belkadiAssociatedMotionDeictic2015,belkadiDeicticDirectionalityAssociated2021,dimmendaalAlloyingEconomyPrinciple2018,dryerAssociatedMotionDirectionals2021,genettiDirectionAssociatedMotion2020,rossCrosslinguisticSurveyAssociated2021}. 

%Given the diverse range of functions of directional extensions, it may be questioned in what sense a directional extension in one language is a comparable entity to a directional extension in another language. For example, a ventive extension in one language might be unique in that it does not share an identical range of functions with a ventive extension in any other language. In order to navigate this complexity, Chapter \ref{ch:distribution} first discusses the distribution of verbal extensions that have a directional function irrespective of any co-expressed function. Chapter \ref{ch:functions} then covers the frequency of various functions that tend to be co-expressed with each semantic type of directional meaning. This is made possible by the database that was constructed for this research in which individual linguistic examples containing a directional extension where coded for the function that the directional expresses in that example, as described in Section \ref{sec:database}. 

\section{Overview}

In Chadic languages, verbal extensions that express direction are frequently described as multifunctional morphemes. The same form may express different meanings depending on the context in which it is used. For an early example of this, in a description of \hyperref[park]{Parkwa}, \citet[315--316]{jarvisPodokoVerbalDirectionals1983} notes that ``directionals may be suffixed to verbs that do not themselves contain the idea of movement. In such cases the directional not only indicates the direction of movement but actually introduces the motion concept.'' Despite the natural inclination to look for motion-related meanings in every use of directional extensions, Jarvis admits that ``sometimes it is difficult to assign any particular meaning to the directional: its use is merely a matter of collocation, just as in English one can speak of ``washing up'' and ``waking up'', without any idea of upward motion.'' 

In light of this multifunctionality, some descriptive work includes very general descriptions of the meaning of directional extensions which are meant to capture the range of uses of the morpheme. \citet[73]{newmanKanakuruLanguage1974} describes a ventive extension in \hyperref[kana]{Kanakuru} as indicating ``that the action took place some distance from, in the direction of, or less often, for the benefit of the speaker.'' \citet[26]{schuhAspectsNgizimSyntax1972} describes the ventive extension in \hyperref[ngiz]{Ngizim} as indicating ``action which takes place in the direction of, or for the benefit of, some person or place of reference (often the speaker or the speaker's location). Sometimes the event itself may have taken place somewhere else, but it ultimately affects the place or person of reference, e.g., `ask' + Ventive = `ask and bring back the answer'.'' The limitations of such general definitions are quickly seen when trying to apply them to particular examples, and even more obvious when any attempt is made to use them to predict how various combinations of a verb and a directional extension will be interpreted (see Chapter \ref{ch:interpretation}).

%\citet[663]{newmanHausaLanguageEncyclopedic2000} on the \hyperref[haus]{Hausa} ventive: ``The ventive ending generally denotes action or movement in the direction of the speaker (or any other pragmatically established deictic center)...'' 

While it is often possible to identify logical or metaphorical connections across many of the functions co-expressed in a direction verbal extension, the range of uses of directional extensions appear to vary unpredictably across languages, even across closely related languages. From a typological perspective, it may be questioned in what sense a directional extension in one language is a comparable entity to a directional extension in another language. For this reason, this chapter avoids discussing each type of directional extension as if it were a natural kind \citep[see][]{haspelmathHowComparativeConcepts2018}. Rather, following the work of \citet{francoisSemanticMapsTypology2008} and \citet{hartmannIdentifyingSemanticRole2016}, the discussion is framed in terms of verbal extensions that co-express multiple meanings of which at least one is a directional meaning (see Chapter \ref{ch:identifying2}, Section \ref{sec:onetomany}). The co-expressed (non-directional) meanings found in the data from Chadic languages are organized into three groups for the purpose of presentation, not necessarily reflecting any typological patterns. 

The first group, discussed in Section \ref{sec:comotion}, gathers together meanings related to motion and location (other than direction). The most common type of motion meaning co-expressed with directional meaning is subsequent motion, as illustrated in (\ref{ex:bole5}).

\ea\label{ex:bole5}
\langinfo{\hyperref[bole]{Bole} verb \textit{gàtt} `get tired' with ventive suffix}{}{\cite[139]{gimbaBoleVerbMorphology2000}} \\
\gll à gàtt-àakóo \\
?? get.tired-\gl{vent} \\
\glt `He will get tired and come.'
\z

The second group, discussed in Section \ref{sec:cogrammatical}, consists of uses that modify the meaning of the predicate in other ways, not related to direction, motion or location. These may be related to tense and aspect, argument structure or other modification of the predicate. One example of this is shown in (\ref{ex:bura0}), where the same verbal extension that elsewhere expresses outward direction is shown here to express an aspectual meaning `thoroughly, completely'. Since these aspectual uses feature prominently in this category, this category will be referred to as aspectual uses, even though the actual range of uses extends beyond aspect. 

\ea\label{ex:bura0}
\langinfo{\hyperref[bura]{Bura} verb \textit{buci} `milk' with outward suffix}{}{\cite[4]{blenchBuraVerbalExtensions2010}} \\
\gll buci-ɓəla	\\
milk-\gl{out}		\\
\glt `to do a thorough job of milking'
\z 

It is also common to find specific combinations of a verb root and a particular verbal extension which have an idiomatic interpretation that does not preserve the core meaning of the verb, as discussed in Section \ref{sec:colexical}. In (\ref{ex:park0}) and (\ref{ex:malg0}), the meaning of the predicate cannot be semantically decomposed into the meaning of its parts. The predicate `answer' cannot be systematically derived by a compositional system that combines the meanings of \textit{mbəɗ} `pour' and the \textit{-arə} `onto' suffix, and likewise for the verb \textit{kə́ɮa} `break' combining with an upward extension \textit{-tə́} to mean `hem, wrap'.

\ea\label{ex:park0}
\langinfo{\hyperref[park]{Parkwa} verb \textit{mbəɗ} `pour' with onto suffix}{}{\cite[320]{jarvisPodokoVerbalDirectionals1983}} \\
\gll Mbəɗ-arə mayá.	\\
pour-\gl{onto} 1\gl{sg}		\\
\glt `I answered.'

\ex\label{ex:malg0}
\langinfo{\hyperref[malg]{Malgwa} verb \textit{kə́ɮa} `break' with upward suffix}{}{\cite[150]{lohrSpracheMalgwaNara2002}} \\
\gll kə́ɮa-tə́	\\
break-\gl{up}		\\
\glt `to hem something (< `to wrap around')'
\z

There are two uses of directional extensions that are discussed in detail in later chapters. One is a common pattern in which a directional extension combines with a verb of acquiring possession (e.g., `take') in order to describe a directed caused accompanied motion event, as in (\ref{ex:kera6}). In Chapter \ref{ch:CAM}, these examples are discussed in a wider context of other forms that express directed caused accompanied motion.

\ea\label{ex:kera6}
\langinfo{\hyperref[kera]{Kera} verb \textit{hàw} ‘take' with ventive extension}{}{\cite[116]{ebertSpracheUndTradition1979}} \\
\gll hàw-dà 	\\
take-\gl{vent}	\\
\glt `Bring it here!' [German: `\textit{Bring es her!}']
\z 

Another context in which it is common to find directional extensions is in making a distinction between the framing of a commercial exchange as a buying or selling event. For example, in \hyperref[gizi]{Giziga}, the verb \textit{hìɗík} can mean `trade, sell, buy', but in its ventive form, \textit{hìɗk-ò}, it explicitly means `buy' \citep[150]{shayGrammarGizigaChadic2021}. In Chapter \ref{ch:buysell}, this use of directional extensions is discussed as one of several strategies for distinguishing buying and selling events across Chadic languages. 

A cursory reading of most descriptions of the multifunctional nature of directional extensions in Chadic languages, considering only the few examples presented to exemplify particular meanings, could give the impression that linguists have a fairly clear grasp on how these verbal extensions are used by speakers of Chadic languages. However, when broadening the scope to examine a larger selection of examples with directional extensions, a different picture emerges. Of the approximately 2,000 unique combinations of a verb and a directional extension in the dataset, in around 1,000 cases the meaning contributed by the verbal extension is unclear or not made explicit in the translation, as in (\ref{ex:haus0}), (\ref{ex:park02}) and (\ref{ex:molo0}). The verbal extensions in each of these examples are elsewhere shown to have directional or subsequent motion meanings. However, in these and other contexts the translation given omits any directional or subsequent motion meaning, and there is no explicit indication in the prose of what the function of the verbal extension is.

\ea\label{ex:haus0}
\langinfo{\hyperref[haus]{Hausa} verb \textit{har̃b} `shoot' with ventive suffix}{}{\cite[662]{newmanHausaLanguageEncyclopedic2000}} \\
\gll  yā har̃b-ō	\\
he shot-\gl{vent}		\\
\glt `He shot (it).'

\ex \label{ex:park02}
\langinfo{\hyperref[park]{Parkwa} verb \textit{uz} `eat' with upward suffix}{}{\cite[94]{jarvisEsquisseGrammaticalePodoko1989}} \\
\gll uz-u mayə́ dafá	\\
eat-\gl{up} 1\gl{sg} boule	\\
\glt `I ate some boule.' [French: `\textit{Je mangeai de la boule}.']

\ex \label{ex:molo0}
\langinfo{\hyperref[molo]{Moloko} verb \textit{wats} `write' with ventive extension}{}{\cite[28]{friesenGrammarMoloko2017}} \\
\gll Hʊrmbʊlɔm à-wats=ala kə ɔkʷɔr aka	\\
God 3\gl{sg}.\gl{pfv}-write=\gl{vent} on stone on	\\
\glt `God wrote them on the stone [tablet].'
\z 

There may be several reasons for this prevalent lack of specificity regarding the function of verbal extensions in the data. It may be an issue of naturalness in translation. A verbal extension might contribute a motion-related meaning that does not translate naturally into the language of analysis (English, French or German) and so could be omitted for that reason. Alternatively, a verbal extension could be contributing a tense-aspect meaning or some other abstract grammatical function which does not have an equivalent in the target translation language, and so is not detectable unless explicitly analyzed by the linguist. A third reason why the function of an extension might be undetectable is that the linguist may incorporate the function of the extension into the gloss of the verb. Imagine, for example, that in (\ref{ex:park0}) the verb \textit{mbəɗ} was glossed `answer' instead of `pour'. In such a case, the function of the verbal extension would be completely obscured. 

It may also be the case that the linguist describing the grammar simply does not yet have a grasp of all of the possible functions of verbal extensions in the language. In a description of a ventive extension in \hyperref[pero]{Pero}, \citet[97]{frajzyngierGrammarPero1989} notes: ``even in simple sentences the functions as described above are not always obvious, or maybe it is the Indo-European background of this linguist that prevents the proper understanding of the choice of ventive rather than non-ventive variants...''

Whatever the reasons for so many examples of verbal extensions with an unclear function, it should be kept in mind that there are many examples that cannot be included in the analysis due to a lack of clarity about the function of the directional extension. Among the examples of directional extensions found in non-motion contexts, only about half of the available data specifies what the function of the extension is. It is unknown whether the other half would exhibit similar patterns if more precise translations were made available, or if the lack of specificity in the unclear cases is due to those extensions having more subtle functions that have not yet been consistently described across Chadic languages. 

Table \ref{tab:DIR_plus} provides an overview of how frequently various meanings are co-expressed with types of directional meaning across Chadic languages. The first column (DIR) repeats from Table \ref{tab:directionals} in Chapter \ref{ch:distribution} how many of the 91 languages in the data have a directional extension of each type of path (ventive, itive, etc.). The second column (+AM/LOC) indicates how many languages show a pattern of co-expressing an associated motion\index{associated motion} or locational meaning in the same verbal extension that expresses a directional meaning. This means that in the 47 languages with a ventive directional extension, that morpheme can also express a motion or location meaning in 40 of those languages. In contrast, in the 19 languages with an itive directional extension, only seven are shown to co-express another motion or location meaning. 

The third column (+Aspect) shows how many languages show evidence of aspectual (and other adverbial) meanings being co-expressed with each directional meaning. The fourth column (+Idiom) counts languages where a directional extension is attested in idiomatic uses. The fifth column (+CAM) tabulates co-expression of caused accompanied motion. The final column (+Unclear) gives the number of languages in which a directional extension is also attested in examples where the meaning contributed by the extension is not discernible in the available data. Note that the issue of an indiscernible function of a verbal extension is widespread across languages and paths of motion.\footnote{In some cases, the lack of any ``unclear'' examples is due to a limited amount of data, rather than an abundance of exceptionally explicit translations.} 

%data in tab:DR_plus is from Python code in Functions 
\begin{table}
	\caption{Number of languages with other meanings co-expressed in a directional extension}
	\label{tab:DIR_plus}
	\begin{tabular}{lrrrrrr} % add l for every additional column or remove as necessary
		\lsptoprule
		{} &  DIR &  +AM/LOC &  +Aspect &  +Idiom &  +CAM &  +Unclear \\
		\midrule
VENT &   47 &                40 &                 6 &             7 &    24 &        30 \\
ITIVE &   19 &                 7 &                 0 &             9 &    11 &        14 \\
INTO &   16 &                 4 &                 2 &             3 &     5 &        11 \\
OUT &   15 &                 3 &                 4 &             7 &    10 &        15 \\
UP &    8 &                 5 &                 1 &             5 &     4 &         8 \\
DOWN &   13 &                 4 &                 1 &             3 &    10 &        11 \\
ONTO &    9 &                 4 &                 4 &             6 &     0 &         7 \\
UNDER &    3 &                 2 &                 0 &             2 &     0 &         3 \\
		\lspbottomrule
%VENT &   47 &                40 &                 6 &             7 &    24 &        30 \\
%dghwOUT overlaps with VENT
%VENT.OUT &    1 &                 1 &                 0 &             1 &     0 &         1 \\ 
%dghwVENT.OUT overlaps with VENT 
%VENT.DOWN &    1 &                 1 &                 0 &             0 &     0 &         0 \\
%ITIVE &   19 &                 7 &                 0 &             9 &    11 &        14 \\
%ITIVE.DOWN &    1 &                 0 &                 0 &             0 &     0 &         0 \\
%INTO &   16 &                 4 &                 2 &             3 &     5 &        11 \\
%OUT &   15 &                 3 &                 4 &             7 &    10 &        15 \\
%UP &    8 &                 5 &                 1 &             5 &     4 &         8 \\
%DOWN &   13 &                 4 &                 2 &             3 &    10 &        11 \\
%ONTO &    9 &                 4 &                 4 &             6 &     0 &         7 \\
%UNDER &    3 &                 2 &                 0 &             2 &     0 &         3 \\
%BUSH &    1 &                 0 &                 1 &             1 &     1 &         1 \\
%HOME &    2 &                 0 &                 0 &             1 &     1 &         2 \\
%SIDE &    1 &                 1 &                 0 &             1 &     0 &         1 \\
%?? &    8 &                 6 &                 0 &             4 &     3 &         7 \\
	\end{tabular}
\end{table}

\section{Motion and location} \label{sec:comotion}

%The co-expression of resultative meaning with directional meaning is not very frequently found in the data on Chadic languages. It is found in 24 morphemes in 16 different languages. 

%data in tab:DR_plustable is from Python code in Functions 
%\begin{table}
%	\caption{Directional extensions with additional functions co-expressed}
%	\label{tab:DR_plusTABLE}
%	\begin{tabular}{llllll} % add l for every additional column or remove as necessary
	%		\lsptoprule
	%		{} &  DIR &  +PriorAM &  +ConcAM &  +Subs(C)AM &  +LOC \\
	%		\midrule
	%VENT &   49 &        4 &       1 &         29 &    5 \\
	%ITIVE &   22 &        0 &       0 &          6 &    1 \\
	%INTO &   17 &        0 &       0 &          1 &    1 \\
	%OUT &   17 &        0 &       0 &          2 &    0 \\
	%UP &   12 &        0 &       1 &          3 &    0 \\
	%DOWN &   16 &        0 &       0 &          3 &    0 \\
	%ONTO &   11 &        0 &       0 &          1 &    1 \\
	%UNDER &    3 &        0 &       0 &          0 &    2 \\
	%VENT.OUT &    1 &        0 &       0 &          1 &    0 \\
	%ITIVE.DOWN &    1 &        0 &       0 &          0 &    0 \\
	%		\lspbottomrule
	%	\end{tabular}
%\end{table}

Recent research has highlighted the cross-linguistic prevalence of another mo\-tion-related meaning found in verbal morphology called associated motion \citep{guillaumeAssociatedMotion2021}. Associated motion is defined in contrast to directional meaning. For the purpose of cross-linguistic comparison, the notion of directional meaning is restricted to cases where the morpheme contributes directional information to an utterance which describes a path of motion. In other words, the fact of motion is expressed somewhere other than the morpheme which expresses direction \citep{guillaumeAssociatedMotionSouth2016}. This is distinguished from the case of associated motion in which a verbal morpheme itself introduces the fact of motion to a construction.\footnote{See \citet{belkadiDeicticDirectionalityAssociated2021} for a proposed unified analysis of directionals and associated motion. Belkadi's approach treats the distinction not as one of grammatical category but as a matter of how the grammars of various languages temporally integrate the motion meaning of the grammaticalized morpheme with the meaning of the predicate. Practically speaking, the distinction still remains, whether it is described as separate functions or at the level of event composition.} For example, in (\ref{ex:Kaytetye1}), from the Pama-Nyungan (Australia) language Kaytetye, the verb \textit{are} ‘look’ is a non-motion verb and no motion is introduced elsewhere in the utterance other than the affix \textit{-yene}, which expresses that the subject changes locations (translocative motion) before the activity predicated by the verb occurs. Since a verbal suffix expresses the fact of motion, this is an instance of associated motion.

\ea\label{ex:Kaytetye1}
\langinfo{Kaytetye (Pama-Nyangan) verb \textit{are} ‘look’ with prior motion affix}{}{\cite[238]{kochAssociatedMotionPamaNyungan2021}} \\
\gll are-nke=lke re, wethapenye are-yene-nke weye-le=pe \\
look-\gl{prs}=then 3\gl{sg}.\gl{erg} thus look-\textsc{go}\&\textsc{do}-\gl{prs} animal-\gl{erg}=\gl{top} \\
\glt `It (euro) looks then, the animal goes and has a look.'
\z 

One parameter along which different types of associated motion\index{associated motion} meaning can be distinguished is whether the form expresses a  motion event prior to the activity of the main verb (as with \textit{-yene} in (\ref{ex:Kaytetye1})), or a  motion event subsequent to the activity of the main verb (as with \textit{-layte} in (\ref{ex:Kaytetye2})) or a motion event concurrent with the activity of the main verb (as with \textit{-larelarre} in (\ref{ex:Kaytetye3})).

\ea\label{ex:Kaytetye2}
\langinfo{Kaytetye (Pama-Nyangan) verb \textit{are} ‘look’ with subsequent motion affix}{}{\cite[241]{kochAssociatedMotionPamaNyungan2021}} \\
\gll well re atyenge atere-le are-layte-nhe \\
well 3\gl{sg}.\gl{erg} 1\gl{sg}.\gl{acc} frightened-\gl{erg} see-\textsc{do}\&\textsc{go}-\gl{pst} \\
\glt `Well, it (bird) saw me and flew off frightened.'

\ex\label{ex:Kaytetye3}
\langinfo{Kaytetye (Pama-Nyangan) verb \textit{are} ‘look’ with concurrent motion affix}{}{\cite[243]{kochAssociatedMotionPamaNyungan2021}} \\
\gll nharte=pe are-larelarre-yayne atye=pe, are-yernalpe-yayne atye \\
then=\gl{top} look-all.along-\gl{ipfv}.\gl{pst} 1\gl{sg}.\gl{erg}=\gl{top} look-coming-\gl{ipfv}.\gl{pst} 1\gl{sg}.\gl{erg} \\
\glt `Then I kept watching all the way along, I watched as we came this way.'
\z 

Table \ref{tab:DR_plusTABLE2} shows the number of languages that have a verbal extension in which directional meaning is co-expressed with a motion meaning or a location meaning. The column DIR shows the number of languages that have a directional extension of each of the most common types. The other columns indicate in how many languages a directional extension of that type co-expresses another meaning: Prior associated motion (+PriorAM), concurrent associated motion (+ConcAM), subsequent associated motion or subsequent caused accompanied motion (+Subs(C)AM), and location (+LOC).

%data in tab:DR_plustable2 is from Python code in Functions 
\begin{table}
	\caption{Directional extensions and motion/location co-expression}
	\label{tab:DR_plusTABLE2}
	\begin{tabular}{lrrrrr} % add l for every additional column or remove as necessary
		\lsptoprule
		{} &  DIR &  +PriorAM &  +ConcAM &  +Subs(C)AM &  +LOC \\
		\midrule
VENT &   47 &        3 &       0 &         30 &    7 \\
ITIVE &   19 &        0 &       0 &          7 &    0 \\
INTO &   16 &        0 &       0 &          1 &    3 \\
OUT &   15 &        0 &       0 &          3 &    0 \\
UP &    8 &        0 &       1 &          4 &    0 \\
DOWN &   13 &        0 &       0 &          4 &    0 \\
ONTO &    9 &        0 &       0 &          1 &    3 \\
UNDER &    3 &        0 &       0 &          0 &    2 \\
		\lspbottomrule
%OUT &   15 &        0 &       0 &          3 &    0 \\
%VENT.OUT &    1 &        0 &       0 &          1 &    0 \\
%VENT.DOWN &    1 &        0 &       0 &          1 &    0 \\
%VENT &   47 &        3 &       0 &         30 &    7 \\
%ITIVE &   19 &        0 &       0 &          7 &    0 \\
%?? &    8 &        0 &       0 &          4 &    2 \\
%ONTO &    9 &        0 &       0 &          1 &    3 \\
%INTO &   16 &        0 &       0 &          1 &    3 \\
%UP &    8 &        0 &       1 &          4 &    0 \\
%DOWN &   13 &        0 &       0 &          4 &    0 \\
%BUSH &    1 &        0 &       0 &          0 &    0 \\
%UNDER &    3 &        0 &       0 &          0 &    2 \\
%HOME &    2 &        0 &       0 &          0 &    0 \\
%ITIVE.DOWN &    1 &        0 &       0 &          0 &    0 \\
%SIDE &    1 &        0 &       0 &          0 &    1 \\
	\end{tabular}
\end{table}

As is clear from Table \ref{tab:DIR_plus}, verbal extensions that express ventive directional meaning are the most likely to co-express another motion or location meaning.\footnote{By binomial test, 40 of 47 (85\%) is significantly higher (p < 0.01) than the overall rate at which directional extensions are shown to co-express associated motion or location.} In 40 of the 47 languages (85\%) that have a verbal extension that expresses ventive directional meaning, the same form is also reported to have another motion or location meaning.\footnote{The \hyperref[baca]{Bacama} and \hyperref[mbud]{Mbudum} ventive extensions co-express both prior and subsequent motion. The \hyperref[muya]{Muyang}, \hyperref[haus]{Hausa}, \hyperref[bole]{Bole} and \hyperref[park]{Parkwa} ventive extensions co-express both subsequent motion and locative. The \hyperref[kana]{Kanakuru} ventive forms co-express prior motion, subsequent motion and locative.} Table \ref{tab:DR_plusTABLE2} reveals that the most commonly found meaning co-expressed with directional meaning is subsequent motion. It is rare to find other types of motion or location meaning co-expressed with directional meaning.

In Chadic languages, many (but not all) verbal extensions that have a directional function also have an associated motion\index{associated motion} function. Rarely (if ever) do verbal extensions express associated motion without co-expressing directional meaning.\footnote{There are three possible exceptions in the dataset. The \hyperref[psik]{Psikye} extension \textit{-ke} expresses ventive subsequent motion. It is not shown to have a ventive directional function in any examples, though the same form in a closely related language does. The \hyperref[glav]{Glavda} extension \textit{-al} is said to mean `separation or apart' \citep[25]{nghagyiyaSketchGlavdaGrammar2011}, but in one example it has a possible itive subsequent motion interpretation: ``stand and walk away''. The \hyperref[gaan]{Ga'anda} extension \textit{-in} `along' is twice given a ventive subsequent motion interpretation, but never a directional interpretation. In addition, the \hyperref[saku]{Sakun} suffixes \textit{-kə́} `ventive' and \textit{-rá} `centrifugal' are not found with directional interpretations, but they are used in directed CAM expressions with `take' which can be viewed as a type of subsequent associated motion.} For this reason, it is appropriate to consider these extensions primarily directional, with associated motion\index{associated motion} as an extended use of directional extensions. This stands in contrast to other language families that exhibit more robust associated motion morphology which might be extended to directional uses \citep{dryerAssociatedMotionDirectionals2021}.\footnote{For example, in a few Atlantic languages associated motion can be identified as the primary function with directional meaning as a secondary function \citep{voisinAssociatedMotionDeictic2021}, and in Choctaw directional morphemes and associated motion morphemes do not overlap in their functions at all \citep{broadwellChoctawDirectionalsSyntax2000}.}

The following subsections examine in more detail the types of motion and location meaning co-expressed with each type of directional meaning. Section \ref{sec:covent} discusses co-expression with ventive directional extensions. These cases exemplify the common pattern of co-expression of direction and subsequent motion, but also include the few languages in which there is a co-expression of direction and prior associated motion. It is of note that in the co-expression of prior associated motion the direction of the prior motion is itive, despite being co-expressed with a ventive directional extension.

Other types of directional meaning are typically only co-expressed with subsequent motion or, in a few cases, a static locative meaning. Co-expression with itive direction is discussed in Section \ref{sec:coitive}. Co-expression with boundary-crossing directions (inward and outward) is discussed in Section \ref{sec:cobounded}. Co-expression with vertical direction is discussed in Section \ref{sec:covert}. This includes the only attested case of concurrent motion co-expressed with directional meaning. Co-expression of motion meanings with `onto' and `under' direction are not discussed due to an insufficient number of examples in the data.

\subsection{Co-expression with ventive directional extensions} \label{sec:covent}

As summarized above in Table \ref{tab:DR_plusTABLE2}, the majority of ventive directional extensions are described as being co-expressive with other motion and location meanings. Table \ref{tab:vent_plusTABLE} shows how many languages provide evidence of each type of co-expressed meaning across the four branches of Chadic languages. Co-expression of subsequent motion with ventive direction is found to be much more common that prior or concurrent motion, as already observed by  \citet[253]{bourdinAllerVenirSomali2005}. There are, however, some exceptional cases of co-expression with prior associated motion\index{associated motion} or locative meaning attested in Biu-Mandara and West Chadic languages.

%Data in table vent_plusTABLE is tabulated in Python file FunctionsBranch
\begin{table}
	\caption{Ventive directional extensions and co-expression by branch}
	\label{tab:vent_plusTABLE}
	\begin{tabular}{lrrrrr} % add l for every additional column or remove as necessary
		\lsptoprule
           {} &  Ventive &  +PriorAM &  +ConcAM &  +Subs(C)AM &  +LOC \\
\midrule
West &       15 &         1 &        0 &          10 &     4 \\
Biu-Mandara & 26 &         2 &        0 &          18 &     3 \\
Masa &        3 &         0 &        0 &           0 &     0 \\
East &        3 &         0 &        0 &           2 &     0 \\
Total &       47 &         3 &        0 &          30 &     7 \\
		\lspbottomrule
	\end{tabular}
\end{table}

\subsubsection{Ventive direction and subsequent motion}

Where subsequent motion is co-expressed with a ventive directional extension, as in (\ref{ex:haus3}) and (\ref{ex:lele3}), the direction of the subsequent motion is also ventive. 

\ea\label{ex:haus3}
\langinfo{\hyperref[haus]{Hausa} verb \textit{shāf} `whitewash' with ventive suffix}{}{\cite[663]{newmanHausaLanguageEncyclopedic2000}} \\
\gll nā	shāf-ō	bangō	 \\
I	whitewash-\gl{vent}	wall \\
\glt `I whitewashed the wall and came back.'

\ex\label{ex:lele3}
\langinfo{\hyperref[lele]{Lele} verb \textit{lèé} ‘eat’ with ventive marker}{}{\cite[194]{frajzyngierGrammarLele2001}} \\
\gll ŋ	lèé	jè	sìí \\
1\gl{sg}	eat	\gl{vent}	meat \\
\glt `I ate the meat [somewhere else and came here].'
\z 

The subsequent motion function of a directional extension may be made less obvious if the same path of motion is also expressed elsewhere in the immediate discourse context.\footnote{The use of verbal morphology to indicate a subsequent motion event that is repeated in the immediate context is a pattern that has been observed in other languages with associated motion morphology (\citeauthor{guillaumeIntroductionAssociatedMotion2021} \citeyear[15]{guillaumeIntroductionAssociatedMotion2021}; \citeauthor{wilkinsSemanticsPragmaticsDiachronic1991} \citeyear[251]{wilkinsSemanticsPragmaticsDiachronic1991}).} For example, in \hyperref[cuvo]{Cuvok}, there are unambiguous examples of the ventive directional co-expressing subsequent motion; however, in (\ref{ex:cuvo02}) there is some ambiguity. If the ventive directional extension on the verb \textit{ndə̀v} `finish' is interpreted as expressing a subsequent motion event, then this meaning is redundant with the following clause. 

\ea\label{ex:cuvo02}
\langinfo{\hyperref[cuvo]{Cuvok} verb \textit{ndə̀v} `finish' with ventive suffix}{}{\cite[305]{dadakGrammaireCuvokLangue2021}} \\
\gll kɛ̀-ndə̀v-ɛ́k ɮɛ́lɛ̀j tá-kà, tá kɛ̀-vɛ̀h-ɛ́k vàrá pápáŋ \\
3\gl{sg}.\gl{sbj}.\gl{pst}-finish-\gl{vent} wealth \gl{asoc}-2\gl{sg}.\gl{poss} then 3\gl{sg}.\gl{sbj}.\gl{pst}-return-\gl{vent} \gl{loc} father.3\gl{sg}.\gl{poss} \\
\glt `(Before coming here) he wasted all his fortune and then returned (here) to his father.' \newline
[French: `\textit{Il a gaspillé toute ta richesse et il est retourné vers son père}.']
\z

The \hyperref[gava]{Gavar} example in (\ref{ex:gava0}) is similar. There are two non-motion verbs, both of which have a ventive directional extension which could potentially be expressing subsequent motion. The translation has just one subsequent motion event, which presumably takes place after the activities predicated by both verbs are completed. In other words, it is not the case that the subject returns after the field is clear, and then goes back out to burn the field before returning a second time. 

\ea\label{ex:gava0}
\langinfo{\hyperref[gava]{Gavar} verb \textit{tsax} `clear' with ventive suffix}{}{\cite[36]{viljoenGavarVerbPhrase2017}} \\
\gll ndə́ tsáx xì-tsax-xà là bə̀zà féŋ xì-feŋ-xà là bə̀zà	 \\
\gl{rel}.\gl{obj} clear 1\gl{incl}.\gl{sbj}-clear-\gl{vent} field \gl{compl} burn 1\gl{incl}.\gl{sbj}-burn-\gl{vent} field \gl{compl}	 \\
\glt `When we clear the field, we burn the field (before returning home).'
\z 

In this data, the category of \textit{subsequent motion} conflates the notionally distinct concepts of subsequent associated motion\index{associated motion} and subsequent caused accompanied motion. This is because there are no examples of a directional extension that is shown to co-express subsequent caused accompanied motion without also co-expressing subsequent associated motion. 

Examples (\ref{ex:haus3}) and (\ref{ex:lele3}) above are considered subsequent associated motion because the added meaning is simply that one of the arguments of the predicate (the subject) is also a figure moving along a path of motion. Caused accompanied motion (CAM) is a label used to describe `bring' events where the figure moving along a path of motion is described as causing another argument to move along the same path at the same time (Chapter \ref{ch:CAM}). In transitive sentences, as in (\ref{ex:haus5}) and (\ref{ex:baca5}), it is sometimes the case that a subsequent motion interpretation includes a caused accompanied motion meaning: the agent causes the patient to move concurrently on the path of subsequent motion. 

\ea\label{ex:haus5}
\langinfo{\hyperref[haus]{Hausa} verb \textit{say} `buy' with ventive suffix}{}{\cite[663]{newmanHausaLanguageEncyclopedic2000}} \\
\gll yā	say-ō	nāmā̀ \\
he	buy-\gl{vent}	meat \\
\glt `He bought some meat and brought it back here.'

\ex\label{ex:baca5}
\langinfo{\hyperref[baca]{Bacama} verb \textit{ɓay} `break' with ventive suffix}{}{\cite[102]{carnochanCategoriesVerbalPiece1970}} \\
\gll Hūn ɓay-a kàdē	\\
1\gl{sg} break-\gl{vent} sticks		\\
\glt `I have broken the sticks and brought them.'
\z 

In the above examples of the \hyperref[haus]{Hausa} ventive form, the interpretation in (\ref{ex:haus3}) is subsequent associated motion,\index{associated motion} while the interpretation in (\ref{ex:haus5}) is subsequent caused accompanied motion. The reason for the different interpretations appears to be pragmatic. The object `wall' in (\ref{ex:haus3}) is not something a person easily carries, and the assumption (based on knowledge of the world) is that the wall was painted in the place where it was meant to stand rather than being moved after being painted. In contrast, the object `meat' in (\ref{ex:haus5}) is presumably portable, and the buying event described is one often done with the intention of transporting the goods purchased. This can be compared with the \hyperref[lele]{Lele} example (\ref{ex:lele3}) where the object is also `meat', but since the predicate `eat' involves consuming the patient, a subsequent CAM interpretation would not be expected (even though technically the meat is transported in some form). In summary, the CAM interpretation of subsequent motion in Chadic languages is pragmatically inferred. 

\subsubsection{Ventive direction and itive prior motion}

In three languages, prior associated motion\index{associated motion} (`go and then...') is co-expressed with ventive directional meaning in the same verbal extension. In all three cases, the direction of prior motion does not match the ventive directional interpretation, but is rather itive, as in (\ref{ex:kana3}), (\ref{ex:mbud14}) and (\ref{ex:baca4}). 

\ea\label{ex:kana3}
\langinfo{\hyperref[kana]{Kanakuru} verb \textit{ɗop} `tie' with ventive suffix}{}{\cite[7]{newmanKanakuruLanguage1974}} \\
\gll a	ɗop-təru \\
\gl{pfv} tie-\gl{vent} \\
\glt `He went and tied it.'

\ex\label{ex:mbud14}
\langinfo{\hyperref[mbud]{Mbudum} verb \textit{sə̀kə̀m} `buy' with ventive suffix}{}{\cite[289]{fobakayyangPhonologieMorphologieSyntaxe2021}} \\
\gll à=sə̀kə̀m-àhà ndə̀rə̀j ə́ŋg lə́mà hwàsàm	 \\
3\gl{sg}.\gl{sbj}=buy-\gl{vent} millet \gl{prep} market Hosom	\\
\glt `He goes and buys millet at the market in Hosom.' \newline
[French: `\textit{Il part achèter le mil au marché de Hosom}.']

\ex\label{ex:baca4}
\langinfo{\hyperref[baca]{Bacama} verb \textit{ɗàw} `cut' with ventive suffix}{}{\cite[93]{carnochanCategoriesVerbalPiece1970}} \\
\gll Ndā ɗàw-á-go Pwèddon kadā	\\
3\gl{sg}.\gl{m} cut-\gl{vent}-\textsc{deprivative} Pweddon tree		\\
\glt `He went and cut down the tree without Pweddon's knowledge.'
\z 

In the case of \hyperref[baca]{Bacama} in (\ref{ex:baca4}), the use of the ventive suffix with non-motion verbs can also result in a subsequent motion interpretation, in which case the direction of the motion is ventive, as in (\ref{ex:baca3}). It is not known what factors condition whether the extension is given a prior or subsequent motion interpretation.

\ea\label{ex:baca3}
\langinfo{\hyperref[baca]{Bacama} verb \textit{ɗàw} `cut' with ventive suffix}{}{\cite[92]{carnochanCategoriesVerbalPiece1970}} \\
\gll Ndā ɗàw-á Victòr kadā	\\
3\gl{sg}.\gl{m} cut-\gl{vent} Victor tree		\\
\glt `He cut down the tree for Victor and returned.'
\z 

Not included in the statistics for prior motion are several cases where the translation indicates round-trip motion (e.g., went, did something, and came back), as in (\ref{ex:zulg3}) and (\ref{ex:mofu5}). 

\ea\label{ex:zulg3}
\langinfo{\hyperref[zulg]{Zulgo} verb \textit{tə̀f} `draw' with ventive suffix}{}{\cite[48]{hallerVerbalComplexZulgo1981}} \\
\gll á-tə̀f-ára yam-á \\
she-drew-\gl{vent} water-here \\
\glt `She went, drew water and brought it here.'

\ex\label{ex:mofu5}
\langinfo{\hyperref[mofu]{Mofu} verb \textit{səp} `search' with ventive suffix}{}{\cite[9]{hollingsworthTransitivityPragmaticsObject1995}} \\
\gll  a səp-ər-wa ɗakw aŋga	\\
3 search-\gl{obj}-\gl{vent} goat his		\\
\glt `(He goes and) looks for his goat (to bring it back).'
\z 

The round-trip motion interpretation is only found co-expressed in ventive directional extensions. In those Chadic languages where a round-trip motion interpretation is found, there are typically other examples showing that a subsequent motion interpretation is also possible (without expressing a round trip), as in (\ref{ex:zulg20}).

\ea\label{ex:zulg20}
\langinfo{\hyperref[zulg]{Zulgo} verb \textit{sə̀kə́m} `buy' with ventive suffix}{}{\cite[47]{hallerVerbalComplexZulgo1981}} \\
\gll á-sə̀kə́m-ára slú i kwàskwà-ya	\\
he-bought-\gl{vent} meat at market-here		\\
\glt `He bought meat at the market and brought it here.'
\z 

Given the pattern of co-expression of subsequent motion and round-trip motion, the attested cases of round-trip motion are treated here as subsequent motion meaning with a context-based inference that an itive path of motion must have taken place before the activity described by the verb in order for the moving figure to be able to make the ventive path of motion subsequent to that activity. Therefore round-trip motion is not treated as a separate semantic category in this analysis. It is considered a pragmatic inference.

The examples of round-trip inference in the context of subsequent motion offer a possible explanation for how ventive direction comes to be co-expressed with itive prior motion. The historical steps would progress from only expressing ventive subsequent motion, to expressing both itive prior and ventive subsequent motion (round trip), to only expressing itive prior motion. However, this path of semantic shift is by no means inevitable. 

In other language families, there are languages where a verbal morpheme co-expresses ventive directional meaning and ventive prior motion (rather than itive prior motion). For example, \citet[173]{mietznerSpatialOrientationNilotic2012} cites the case of Päri (Nilotic), shown in (\ref{ex:pari}). Notice that the Päri example in (\ref{ex:pari1}) shows that this language also allows the itive directional marker to co-express itive prior motion. \citet[79--80]{christinelamarreDeicticDirectionalsRevisited2022} document a similar pattern in Mandarin (Sino-Tibetan) and Sereer (Niger-Congo).

\ea\label{ex:pari}
\ea 
\langinfo{Päri (Nilotic) verb \textit{ŋùnn} `cut' with ventive suffix}{}{\cite[88]{andersenConsonantAlternationVerbal1988}} \\
\gll  ùbúr á-ŋùnn-ò \\
Ubur \gl{compl}-cut.\gl{vent}.\gl{antip}-\gl{intr} \\
\glt `Ubur came to cut.'

\ex  \label{ex:pari1}
\langinfo{Päri (Nilotic) verb \textit{ŋùnn} `cut' with itive suffix}{}{\cite[88]{andersenConsonantAlternationVerbal1988}} \\
\gll  ùbúr á-ŋùt-ò \\
Ubur \gl{compl}-cut.\gl{itive}.\gl{antip}-\gl{intr} \\
\glt `Ubur went to cut.'
\z 
\z

\subsubsection{Ventive direction and locative} \label{sec:ventloc}

There are seven languages in which a verbal extension that expresses ventive directional meaning is attested as co-expressing a locative meaning. As shown in Table \ref{tab:vent_plusTABLE}, three of those languages are Biu-Mandara languages (\hyperref[mina]{Mina}, \hyperref[muya]{Muyang} and \hyperref[mada]{Mada}) and four are West Chadic languages (\hyperref[kana]{Kanakuru}, \hyperref[bole]{Bole}, \hyperref[shaa]{Ron-Scha} and \hyperref[haus]{Hausa} in at least one example). In five of these languages, the locative meaning co-expressed with a ventive directional extension has an altrilocative (`there') interpretation, as in (\ref{ex:mada2}), (\ref{ex:mina7}) and (\ref{ex:kana8}).\footnote{The \hyperref[kana]{Kanakuru} ventive form is also shown to have an itive prior motion interpretation in some cases, as in (\ref{ex:kana3}). It is not known in what contexts the different interpretations are found.}

\ea\label{ex:mada2}
\langinfo{\hyperref[mada]{Mada} verb \textit{zlèc} `hit' with ventive suffix}{}{\cite[46]{ernstkurdiTenseAspectMood2016}} \\
\gll  zlèc-ēré \\
hit.\gl{imp}-\gl{vent}		\\
\glt `Hit (in another place)!'

\ex\label{ex:mina7}
\langinfo{\hyperref[mina]{Mina} verb \textit{màr} `shepherd' with ventive suffix}{}{\cite[175]{frajzyngierGrammarMina2005}} \\
\gll í	ndí	màr-áhà		\\
3\gl{pl}	\gl{hab}	shepherd-\gl{vent}	 \\
\glt `They have the habit of pasturing somewhere else.'

\ex\label{ex:kana8}
\langinfo{\hyperref[kana]{Kanakuru} verb \textit{àl} `see' with ventive suffix and pronominal suffix}{}{\cite[74]{newmanKanakuruLanguage1974}} \\
\gll mə̀	àl-də̀	ré \\
1\gl{pl}	see-\gl{vent}	3\gl{sg}.\gl{f}.\gl{obj} \\
\glt `We saw her (at a distance).'
\z 

In two languages, \hyperref[bole]{Bole} and \hyperref[shaa]{Ron-Scha}, the verbal extension that co-expresess ventive directional meaning and locative meaning has a cislocative (`here') interpretation in its locative use, as in (\ref{ex:bole20}) and (\ref{ex:bole2}). 

\ea\label{ex:bole20}
\langinfo{\hyperref[bole]{Bole} verb \textit{yòrú} `stop' with ventive suffix}{}{\cite[138]{gimbaBoleVerbMorphology2000}} \\
\gll yòrú-n nzònó	\\ 
stop-\gl{vent} yesterday		\\
\glt `He stopped here yesterday.'

\ex\label{ex:bole2}
\langinfo{\hyperref[bole]{Bole} verb \textit{ngòrú} `tie' with ventive suffix}{}{\cite[142]{gimbaBoleVerbMorphology2000}} \\
\gll íshí	ngòrú-ttùu-yì \\
he	tie-\gl{vent}-\gl{obj} \\
\glt `that he tie (it) here'
\z 

There is, however, some ambiguity concerning the cislocative interpretation in \hyperref[bole]{Bole}. The verb \textit{yòrú} `stop' with a ventive suffix is given a cislocative intepretation in (\ref{ex:bole20}) and (\ref{ex:bole2}), but there is a third example, in (\ref{ex:bole111}), where it is given an altrilocative translation. The same verbal extension is more commonly shown with a subsequent motion interpretation with non-motion verbs, as in (\ref{ex:bole1}). Note that the same verb appears in (\ref{ex:bole2}) and (\ref{ex:bole1}) in a similar morphosyntactic context. It is not clear what factors condition the interpretation of this verbal extension. 

\ea\label{ex:bole111}
\langinfo{\hyperref[bole]{Bole} verb \textit{'yòr} `stop' with ventive suffix}{}{\cite[139]{gimbaBoleVerbMorphology2000}} \\
\gll à 'yòr-aàkóo \\
? stop-\gl{vent}		\\
\glt `he will stop (there)'	

\ex\label{ex:bole1}
\langinfo{\hyperref[bole]{Bole} verb \textit{ngòrú} `tie' with ventive suffix}{}{\cite[143]{gimbaBoleVerbMorphology2000}} \\
\gll íshí	ngòrú-ttú	dóoshó \\
he	tie-\gl{vent}	horse \\
\glt `that he tie a horse and bring here'
\z 

In summary, while there are some clear examples of ventive extensions with a locative function, the available description and translations are often imprecise, creating some ambiguity around the details of what types of locative functions are co-expressed with ventive direction. 

\subsection{Co-expression with itive directional extensions} \label{sec:coitive}

As discussed in Chapter \ref{ch:distribution}, itive directional extensions are less common than ventive directional extensions. They are primarily limited to Biu-Mandara and Masa languages with one exception. Table \ref{tab:itive_plusTABLE} shows that, of the relevant types of meaning, only subsequent motion is found to be co-expressed with itive directional meaning in the same verbal extension, and only in Biu-Mandara languages.

%Data in table vent_plusTABLE is tabulated in Python file FunctionsBranch
\begin{table}
	\caption{Itive directional extensions and co-expression by branch}
	\label{tab:itive_plusTABLE}
	\begin{tabular}{lrrrrr} % add l for every additional column or remove as necessary
		\lsptoprule
	{} &  Itive &  +PriorAM &  +ConcAM &  +Subs(C)AM &  +LOC \\
		\midrule
West &      0 &         0 &        0 &           0 &     0 \\
Biu-Mandara &     16 &         0 &        0 &           7 &     0 \\
Masa &      2 &         0 &        0 &           0 &     0 \\
East &      1 &         0 &        0 &           0 &     0 \\
\midrule
Total &     19 &         0 &        0 &           7 &     0 \\
		\lspbottomrule
	\end{tabular}
\end{table}

The seven Biu-Mandara languages in which itive directional meaning is co-expressed with subsequent motion are \hyperref[buwa]{Buwal}, \hyperref[cuvo]{Cuvok}, \hyperref[glav]{Glavda}, \hyperref[psik]{Psikye}, \hyperref[zulg]{Zulgo-Gemzek} and \hyperref[bana]{Bana}. In all seven languages, the direction of the subsequent motion interpretation is itive, as in (\ref{ex:cuvo9}) and (\ref{ex:psik2}). 

\ea\label{ex:cuvo9}
\langinfo{\hyperref[cuvo]{Cuvok} verb \textit{hə̀v} `cultivate' with itive suffix}{}{\cite[308]{dadakGrammaireCuvokLangue2021}} \\
\gll  hə̀v-áɗ \\
cultivate-\gl{itive} \\
\glt `Cultivate and go there.'

\ex\label{ex:psik2}
\langinfo{\hyperref[psik]{Psikye} itive form of \textit{tle} `cut'}{}{\cite[113]{smithKapsikiLanguage1969}} \\
\gll  tle-gu le xwá	\\
cut-\gl{itive} \gl{ins} knife	\\
\glt `Cut it with the knife and take it away!'
\z

In the case of \hyperref[zulg]{Zulgo-Gemzek}, there are two examples of subsequent motion co-expression. In (\ref{ex:zulg6}) the direction is itive. In (\ref{ex:zulg7}) the direction is ventive. It is unclear why this is the case. 

\ea\label{ex:zulg6}
\langinfo{\hyperref[zulg]{Zulgo} verb \textit{zə́m} `eat' with itive suffix}{}{\cite[47]{hallerVerbalComplexZulgo1981}} \\
\gll á-zə́m-áha daf \\
he-ate-\gl{itive} fufu \\
\glt `He ate fufu and went there.'

\ex\label{ex:zulg7}
\langinfo{\hyperref[zulg]{Zulgo} verb \textit{vak} `roast' with itive suffix}{}{\cite[47]{hallerVerbalComplexZulgo1981}} \\
\gll á-vak-áha mendzíkwir \\
she-roasted-\gl{itive} chicken \\
\glt `She roasted a chicken and brought it here.'
\z 

\subsection{Co-expression with boundary-crossing directional extensions} \label{sec:cobounded}

Verbal extensions expressing boundary-crossing directional meaning (inward or outward) are only found in Biu-Mandara languages. Sixteen languages have inward directional extensions and 15 have outward directional extensions, as described in Chapter \ref{ch:distribution}, Section \ref{sec:bounded2}. In most cases, no motion- or location-related meanings are co-expressed with these boundary-crossing directional meanings. 

In regard to inward directional meaning, there is just one language where subsequent motion is shown to be co-expressed in the same verbal extension. \citet{jarvisPodokoVerbalDirectionals1983} gives several examples of the verb \textit{ngwad} `attach, bind' with a directional suffix that has a subsequent motion meaning, as in (\ref{ex:park6}). 

\ea\label{ex:park6}
\langinfo{\hyperref[park]{Parkwa} verb \textit{ngwad} `bind' with inward suffix}{}{\cite[315]{jarvisPodokoVerbalDirectionals1983}} \\
\gll Ngwaɗ-akwa mayə́ paní daká.	\\
bind-\gl{into} 1\gl{sg} stalk into.house	\\
\glt `I tied up the stalks and took them into the house.'
\z

However, one example of the same verb with a directional suffix does not have a subsequent motion interpretation, shown in (\ref{ex:park7}). In this case, the locative phrase is a container, so it is not interpreted as the destination of the agent's path of motion, but rather the patient is interpreted as the figure on the path of motion into the sack (`put the corn in the sack'). Logically, the binding event (`tied it up') predicated by the verb must take place after the motion event, so this cannot be a case of subsequent motion. This suggests that the subsequent motion interpretation is pragmatically determined. 

\ea\label{ex:park7}
\langinfo{\hyperref[park]{Parkwa} verb \textit{ngwad} `bind' with inward suffix}{}{\cite[315]{jarvisPodokoVerbalDirectionals1983}} \\
\gll ʸNgwaɗ-akwa mayə́ həyá da bəhwa\\
bind-\gl{into} 1\gl{sg} corn \gl{prep} sack \\
\glt `I put the corn in the sack and tied it up.'
\z

In three languages (\hyperref[glav]{Glavda}, \hyperref[mata]{Matal} and \hyperref[park]{Parkwa}) a verbal extension that expresses inward directional motion may co-express a locative meaning, as in (\ref{ex:glav144}). 

\ea\label{ex:glav144}
\langinfo{\hyperref[glav]{Glavda} verb \textit{ndz} `remain' with inward extension}{}{\cite{owensGlavdaRuralUrban2011}} \\
\gll ndz-əm	\\ 
remain-\gl{into}	\\
\glt `stay in'
\z 

In regard to outward directional meaning, there are three languages where this meaning may be co-expressed with subsequent outward motion: \hyperref[dghw]{Dghwede}, \hyperref[lama]{Lamang} and \hyperref[park]{Parkwa}, as shown in (\ref{ex:dghw}) and (\ref{ex:park}).

\ea\label{ex:dghw}
\langinfo{\hyperref[dghw]{Dghwede} verb \textit{tə́gə̀} `cook' with outward extension}{}{\cite[6]{frickVerbalSystemDghwede1978}} \\
\gll  kâ' tə́gə̀-díle kfì	\\
\gl{narr} cook-\gl{vent}.\gl{out} food	\\
\glt `She cooked food and brought it out.'

\ex\label{ex:park}
\langinfo{\hyperref[park]{Parkwa} verb \textit{ndwaɗ} `bind' with outward extension}{}{\cite[315]{jarvisPodokoVerbalDirectionals1983}} \\
\gll Ndwaɗ-əl-á mayə́ həyá	\\
bind-??-\gl{out} 1\gl{sg} corn		\\
\glt `I bound the corn and took it out.'
\z 

\hyperref[park]{Parkwa} stands out as the only language where both inward and outward extensions are found co-expressing another motion or locative meaning. Otherwise, this pattern appears to be rare (or rarely documented) among Chadic languages that have boundary-crossing directional extensions.

\subsection{Co-expression with vertical directional extensions} \label{sec:covert}

As discussed in Chapter \ref{ch:distribution}, Section \ref{sec:vertical2}, vertical directional extensions are only found in Biu-Mandara languages. Eight languages have an upward directional extension. Thirteen have a downward directional extension. As seen in Table \ref{tab:DR_plusTABLE2} above, in four languages upward direction is co-expressed with subsequent upward motion, and in four languages downward direction is co-expressed with subsequent downward motion. 

There was just one example found, in (\ref{ex:dghw5}), where a vertical directional extension is shown to co-express concurrent motion. In this case, the verb is a non-motion verb, and the translation indicates that the activity predicated by this verb continues while the subject concurrently moves along an upward path. Since (\ref{ex:dghw5}) is the only example of concurrent motion given in \hyperref[dghw]{Dghwede}, it should only tentatively be assumed that directional meaning and concurrent motion are co-expressed by this verbal extension until it is possibly to verify the analysis of Dghwede with further research. 

\ea\label{ex:dghw5}
\langinfo{\hyperref[dghw]{Dghwede} verb \textit{táwə̀} `weep' with upward suffix}{}{\cite[6]{frickVerbalSystemDghwede1978}} \\
\gll kâ' táwə̀-gwá táwà dà lùwá \\
\gl{narr} weep-\gl{itive}.\gl{up} \gl{cont} to home \\
\glt `He was crying all the way up home.'
\z

In one other example in the \hyperref[dghw]{Dghwede} data, upward subsequent motion is co-expressed with upward directional meaning, as in (\ref{ex:dghw50}).\footnote{The difference in the forms \textit{-gwá} and \textit{-də̀gwá} appears to correlate with whether the patient is a figure on the path of motion (see Chapter \ref{ch:identifying2}, Section \ref{sec:allomorphy}).} 

\ea\label{ex:dghw50}
\langinfo{\hyperref[dghw]{Dghwede} verb \textit{gwə́yə̀} `search' with upward suffix}{}{\cite[22]{frickVerbalSystemDghwede1978}} \\
\gll kâ' gwə́yə̀-də̀gwá də̀ge wàyá	\\
\gl{narr} search-\gl{vent}.\gl{up} ?? ??		\\
\glt `He looked for food (and brought it) up.'
\z 

In \hyperref[glav]{Glavda}, \hyperref[psik]{Psikye} and \hyperref[park]{Parkwa}, there is likewise a single example of upward subsequent motion in each language, as in (\ref{ex:psik17}) and (\ref{ex:park01}).\footnote{Smith translates the upward form in (\ref{ex:psik17}) as downward subsequent motion: ``...and brought it down.'' According to Psikye speaker Zra Kodji, Smith's translation is an error and the correct translation should involve upward motion.}

\ea\label{ex:psik17}
\langinfo{\hyperref[psik]{Psikye} upward form of \textit{l} `dig'}{}{\cite[120]{smithKapsikiLanguage1969}} \\
\gll a 'yá ké-l-ák-ate \\
\textsc{active} 1\gl{sg} \gl{ptcp}-dig-??-\gl{up}	\\
\glt `I dug (it) and brought it up.'

\ex\label{ex:park01}
\langinfo{\hyperref[park]{Parkwa} upward form of \textit{ʸngwaɗ} `bind'}{}{\cite[315]{jarvisPodokoVerbalDirectionals1983}} \\
\gll ʸngwaɗ-əl-u mayə́ həyá	\\
bind-??-\gl{up} 1\gl{sg} corn		\\
\glt `I bound the corn and took it up.'
\z

The four languages in which downward subsequent motion is attested are \hyperref[dghw]{Dghwede}, \hyperref[glav]{Glavda}, \hyperref[park]{Parkwa} and \hyperref[gude]{Gude}, as in (\ref{ex:dghw51}), (\ref{ex:gude15}) and (\ref{ex:park5}). Again, the data demonstrating that these verbal extensions can express subsequent downward motion are limited. There are just two examples in Parkwa and Glavda and one each in Dghwede and Gude. 

\ea\label{ex:dghw51}
\langinfo{\hyperref[dghw]{Dghwede} verb \textit{càcə̀} `chase' with upward suffix}{}{\cite[22]{frickVerbalSystemDghwede1978}} \\
\gll ká' mdà càcə̀-dáyà	\\
 \gl{narr} ?? chase-\gl{vent}.\gl{down}	\\
\glt `We chased them (and brought them) down.'

\ex\label{ex:gude15}
\langinfo{\hyperref[gude]{Gude} verb \textit{'anyə} `tie' with downward suffix }{}{\cite[362]{schuhChadicCornucopia2017}} \\
\gll  nyi 'anyə-pà \\
I tie-\gl{down}	\\
\glt `I tied (it) and laid (it) on the ground.'

\ex\label{ex:park5}
\langinfo{\hyperref[park]{Parkwa} verb \textit{t} `cook' with downward suffix}{}{\cite[322]{jarvisPodokoVerbalDirectionals1983}} \\
\gll t-ad-a mayə dafá	\\
cook-??-\gl{down} 1\gl{sg} fufu	 \\
\glt `I cooked some fufu and took it down.'
\z 

\section{Aspectual and adverbial meanings} \label{sec:cogrammatical}
%numbers in this section are from the Python file Functions 

The label \textit{aspectual} is used here in a very general sense to refer to a set of aspectual, adverbial or morphosyntactic functions of verbal extensions other than the motion and locative functions described above. The sections below give the examples of a set of meanings that are most clearly illustrated in the available data. Table \ref{tab:DIR_gram} gives a quantitative overview of how many languages show evidence of one of these co-expressed functions. These include a completive or totality function (COMPL) and other tense and aspect meanings (TAM, Section \ref{sec:tense}), marking a benefactive (BEN, Section \ref{sec:beneficiary}), additive meaning (ADD, Section \ref{sec:additive}) and marking the object as affected by the action such that it splits into parts (SPLIT, Section \ref{sec:objectsplit}).\footnote{Another function of directional extensions occasionally discussed in the literature is increasing or decreasing valency. Without a robust analysis of argument structure patterns it is difficult to discern whether an alternation is necessarily licensed by a verbal extension or if a verb root is ambitransitive (or labile), so there is generally not enough data available to illustrate how this may work in a particular language or to analyze from a cross-linguistic comparative perspective.}

%data in tab:DIR_gram from python script in Functions
\begin{table}
	\caption{Directional extensions and aspectual/adverbial uses}
	\label{tab:DIR_gram}
	\begin{tabular}{lrrrrrr} % add l for every additional column or remove as necessary
		\lsptoprule
     Path &  DIR &  COMPL &  TAM &  BEN &  ADD &  SPLIT \\
\midrule
VENT &   47 &      0 &    1 &    5 &    0 &      0 \\
ITIVE &   19 &      0 &    0 &    0 &    0 &      0 \\
INTO &   16 &      0 &    0 &    0 &    0 &      2 \\
OUT &   15 &      3 &    0 &    1 &    0 &      0 \\
UP &    8 &      1 &    0 &    0 &    0 &      0 \\
DOWN &   13 &      1 &    0 &    0 &    0 &      0 \\
ONTO &    9 &      0 &    0 &    0 &    4 &      0 \\
BUSH &    1 &      1 &    0 &    0 &    0 &      0 \\
		\lspbottomrule
	\end{tabular}
\end{table}

\subsection{Completive aspect} \label{sec:tense}

One meaning co-expressed by directional extensions in Chadic languages relates to an action being done completely or thoroughly, as in (\ref{ex:huba0}). This is attested in four languages.\footnote{In addition, the itive ``particle'' in \hyperref[tera]{Tera} and the bushward/itive adverb in \hyperref[lele]{Lele} (not considered verbal extensions) can also co-express completeness. However, itive directional extensions are not shown to co-express grammaticalized functions in any Chadic languages.} The main direction correlated with completeness is outward direction, as in (\ref{ex:huba0}). Co-expression of completeness with outward motion is attested in \hyperref[huba]{Huba}, \hyperref[bura]{Bura} and \hyperref[gude]{Gude}. The Gude verbal extension \textit{-gi} can express either outward or upward direction, and so is counted twice in Table \ref{tab:DIR_gram}. 

\ea\label{ex:huba0}
\langinfo{\hyperref[huba]{Huba} verb \textit{kwàs} `eat' with outward extension}{}{\cite[113]{muazuGrammarKilbaLanguage2009}} \\
\gll kwàs-bìyà	\\
eat-\gl{out}		\\
\glt `You eat all.'
\z 

There is also one example in Huba where a similar meaning is co-expressed with downward motion, shown in (\ref{ex:huba01}), so Huba is also counted twice in Table \ref{tab:DIR_gram}. In studies of \hyperref[bura]{Bura}, \citet[274]{hoffmannUntersuchungenZurStruktur1955} and \citet[3]{blenchBuraVerbalExtensions2010} compare the use of the outward directional extension to mean `thoroughly, completely' to phrasal verbs in English (e.g., `clean out'). 

\ea\label{ex:huba01}
\langinfo{\hyperref[huba]{Huba} verb \textit{wùr} `finish' with downward extension}{}{\cite[108]{muazuGrammarKilbaLanguage2009}} \\
\gll wùr-yà	\\
finish-\gl{down}	\\
\glt `to finish completely'
\z 

In \hyperref[psik]{Psikye}, the bushward directional extension can be used with non-motion verbs, in which case it is said to mean that ``the action is finished (emphatically)'' \citep[114, 163]{smithKapsikiLanguage1969}, as in (\ref{ex:psik20b}). The bushward direction could be considered a more specific case of outward direction, namely, moving out from an inhabited location. 

\ea\label{ex:psik20b}
\langinfo{\hyperref[psik]{Psikye} bushward form of \textit{séxwé} `wipe'}{}{\cite[163]{smithKapsikiLanguage1969}} \\
\gll 'a ke-séxwé-mte xwá	\\
\textsc{active} \gl{ptcp}-wipe-\gl{bush} knife	\\
\glt `(He) wiped the knife dry.'
\z 

While the co-expression of direction and completeness meaning is only explicitly attested in examples from four languages, prose descriptions of other languages refer to similar interpretations of directional extensions. In \hyperref[psik]{Psikye}, \citet[165]{smithKapsikiLanguage1969} mentions that a putative ventive suffix \textit{-ke} is ``often used with no discernible meaning other than to mark an action as completed.'' \citet[186]{newmanCaseGrammarGa1971} describes the \hyperref[gaan]{Ga'anda} particle \textit{xa} as having two meanings: ``The first is associated with a downward direction of the action... The second meaning denotes that the action is both well done and completely done...'' \citet[216--220]{melisDescriptionMasaTchad1999} states that ventive and itive markers have various aspectual functions mostly relating to viewing the activity as completed. In a description of Sakun, \citet[121]{thomasGrammarSakunSukur2014} mentions that ``Verbal extensions tend to exhibit a correlation with the telicity of events, or the viewing of the event as bounded or completed'' though this is said to be an ``indirect'' effect. 

Among other families of Afroasiatic languages, completive aspect of some type is also associated with ventive directional markers, including Berber languages \citep{mettouchiGrammaticalizationDirectionalClitics2011} and Semitic languages \citep{sjorsMotionVoiceMood2023}. It is interesting to note, however, that in Chadic languages it is not necessarily the ventive directional extension that takes on this function. If the completive interpretation is an extension of the goal-oriented nature of most directional semantics, then it is logical that other directions besides ventive could be extended to have a completive function.

Another type of TAM meaning is found co-expressed in \hyperref[buwa]{Buwal} directional extensions. Buwal is unique in having two ventive directional extensions (Chapter \ref{ch:identifying2}, Section \ref{sec:multiple}). \citet[375--376]{viljoenGrammaticalDescriptionBuwal2013} shows that, at least in some contexts, the two extensions can have contrastive aspectual meanings.\footnote{Viljoen describes the contrast in (\ref{ex:buwa03}) and (\ref{ex:buwa04}) as occurring in the context of past time reference. Note that these same extensions can also express subsequent motion.} In (\ref{ex:buwa03}) and (\ref{ex:buwa04}) the same verb \textit{mār} `begin' is used to describe the same event in history. However, in (\ref{ex:buwa03}) the adverbial phrase `on that day' frames the event as a punctual telic event (completed on that day), while in (\ref{ex:buwa04}) the adverbial phrase `in 1400' encourages an atelic interpretation. Correspondingly, in (\ref{ex:buwa03}) ``the distal [ventive] suffix indicates the situation is now finished'' and in (\ref{ex:buwa04}) ``the proximal [ventive suffix] indicates that it may be still ongoing'' \citep[376]{viljoenGrammaticalDescriptionBuwal2013}.

\ea\label{ex:buwa03}
\langinfo{\hyperref[buwa]{Buwal} verb \textit{mār} `begin' with distal ventive extension}{}{\cite[376]{viljoenGrammaticalDescriptionBuwal2013}} \\
\gll xājāk bwāl ā-mār-xā á pès wēsé	\\
country Buwal \gl{sbj}.3\gl{sg}-begin-\gl{vent}.\gl{dist} \gl{prep} day \gl{dem}.\gl{dist}		\\
\glt `The Buwal country began on that day. (That day is over.)'

\ex\label{ex:buwa04}
\langinfo{\hyperref[buwa]{Buwal} verb \textit{mār} `begin' with proximal ventive extension}{}{\cite[376]{viljoenGrammaticalDescriptionBuwal2013}} \\
\gll xājāk bwāl ā-mār-ā á tā blàkʷ téŋɡʷlèŋ á témérè nfáɗ	\\
country Buwal \gl{sbj}.3\gl{sg}-begin-\gl{vent}.\gl{prox} \gl{prep} on thousand one \gl{prep} hundred four		\\
\glt `The Buwal country began in 1400. (The Buwal country still continues to this day.)'
\z 

\subsection{Benefactive} \label{sec:beneficiary}

Another non-motion function co-expres\-sed with ventive direction in several languages is marking a benefactive. A benefactive function is shown to be co-expres\-sed with ventive direction in five languages. In one language, \hyperref[malg]{Malgwa}, a benefactive function is shown to be co-expressed with outward direction, as in (\ref{ex:malg01}).

\ea\label{ex:malg01}
\langinfo{\hyperref[malg]{Malgwa} verb \textit{ɠyə́ga} `stomp' with outward suffix}{}{\cite[152]{lohrSpracheMalgwaNara2002}} \\
\gll ɠyə́ga-sə́	\\
stomp-\gl{out}		\\
\glt `to stomp for someone'
\z 

In some examples, the benefactive meaning correlates with the presence of a pronominal suffix. In \hyperref[mofu]{Mofu-Gudur}, \citet[8, 10]{hollingsworthTransitivityPragmaticsObject1995} describes a benefactive meaning when the ventive directional extension co-occurs with a first- or second-person object suffix (compare \ref{ex:mofu3} and \ref{ex:mofu4}).

\ea\label{ex:mofu3}
\langinfo{\hyperref[mofu]{Mofu} verb \textit{jaw} `tie up' with 2nd-person object suffix}{}{\cite[10]{hollingsworthTransitivityPragmaticsObject1995}} \\
\gll ya jaw-ka. \\
1 tie.up-2\gl{sbj}.\gl{obj}	\\
\glt `I tie you up.'

\ex\label{ex:mofu4}
\langinfo{\hyperref[mofu]{Mofu} verb \textit{jaw} `tie up' with 2nd-person object suffix and ventive suffix}{}{\cite[10]{hollingsworthTransitivityPragmaticsObject1995}} \\
\gll ya jaw-ka-wa \\
1 tie.up-2\gl{sbj}.\gl{obj}-\gl{vent} \\
\glt `I tie for you.'
\z 

The \hyperref[kana]{Kanakuru} dative pronominal in (\ref{ex:kana7}) has a malefactive or source interpretation, but in (\ref{ex:kana6}) the pronominal suffix is followed by a ventive directional extension and is interpreted as benefactive.

\ea\label{ex:kana7}
\langinfo{\hyperref[kana]{Kanakuru} verb \textit{shit} `steal' with pronominal suffix}{}{\cite[73]{newmanKanakuruLanguage1974}} \\
\gll  à	shin-no	dok \\
\gl{pfv}	steal-1\gl{sg}.\gl{dat}	horse \\
\glt `He stole a horse from me.'

\ex\label{ex:kana6}
\langinfo{\hyperref[kana]{Kanakuru} verb \textit{shit} `steal' with ventive suffix and pronominal suffix}{}{\cite[73]{newmanKanakuruLanguage1974}} \\
\gll à	shit-tə-no	dok \\
\gl{pfv}	steal-\gl{vent}-1\gl{sg}.\gl{dat}	horse \\
\glt `He stole a horse for me.'
\z 

Beneficiary interpretations of ventive directional extensions are also found elsewhere in the Afroasitic phylum, including Berber languages \citep{belkadiDeicticDirectionalitySpace2015} and Semitic languages \citep{sjorsMotionVoiceMood2023}. 

\subsection{Additive} \label{sec:additive}

In four of the nine languages with an `onto' directional extension (Chapter \ref{ch:distribution}, Section \ref{sec:onto2}) there is evidence of a co-expressed additive function, often translated as `in addition' or `again', as in (\ref{ex:marg0}), (\ref{ex:bura01}) and (\ref{ex:park03}). 

\ea\label{ex:marg0}
\langinfo{\hyperref[marg]{Margi} verb \textit{ci} `tell' with `onto' suffix}{}{\cite[134]{hoffmannGrammarMargiLanguage1963}} \\
\gll ci-ŋgə́ri	\\
tell-\gl{onto}	\\
\glt `to tell people something again'

\ex\label{ex:bura01}
\langinfo{\hyperref[bura]{Bura} verb \textit{bwa} `boil' with `onto' suffix}{}{\cite[290]{hoffmannUntersuchungenZurStruktur1955}} \\
\gll bwa-nkir	\\
boil-\gl{onto}	\\
\glt `boil again' [German: `\textit{nochmals kochen}']

\ex\label{ex:park03}
\langinfo{\hyperref[park]{Parkwa} verb \textit{səkw} `buy' with `onto' suffix}{}{\cite[94]{jarvisEsquisseGrammaticalePodoko1989}} \\
\gll  səkw-arə mayə́ burə akə́ sləɓa	\\
buy-\gl{onto} 1\gl{sg} salt \gl{prep} meat	\\
\glt `I bought salt in addition to meat.' \newline [French: `\textit{J'achetai du sel en plus de la viande.}']
\z 

Of the four languages where this meaning is attested in the data, three (\hyperref[bura]{Bura}, \hyperref[marg]{Margi} and \hyperref[huba]{Huba}) are closely related members of the Bura-Margi subgroup in the Margi-Mandara-Mofu group of Biu-Mandara. The other, \hyperref[park]{Parkwa}, is a Mandaraic language, but also in the Margi-Mandara-Mofu group.

\subsection{Object split in parts} \label{sec:objectsplit}

In two languages, \hyperref[marg]{Margi} and \hyperref[malg]{Malgwa}, the inward verbal extension is shown to co-express a meaning related to how the object is affected by the predicate, shown in (\ref{ex:marg01}) and (\ref{ex:malg18}). According to \citet[149]{hoffmannGrammarMargiLanguage1963}, ``The derivatives in \textit{-wá} mostly indicate that the action is done in the direction `into' something... The suffix \textit{-wá} is also frequently used to indicate that the object is divided into (two or more) parts.''

\ea\label{ex:marg01}
\langinfo{\hyperref[marg]{Margi} verb \textit{ɓàts} `break' with inward suffix}{}{\cite[148]{hoffmannGrammarMargiLanguage1963}} \\
\gll ɓàts-wá	\\
break-\gl{into}	\\
\glt `to break in two pieces'

\ex\label{ex:malg18}
\langinfo{\hyperref[malg]{Malgwa} verb \textit{kya} `break' with inward suffix }{}{\cite[162]{lohrSpracheMalgwaNara2002}} \\
\gll kya-an-ə́m	\\
break-\gl{tr}-\gl{into}	\\
\glt `to break in two' [German: `\textit{in zwei Teile brechen}']
\z

\section{Idiomatic verb-extension combinations} \label{sec:colexical}

%Numbers in this paragraph from Python code in Functions. See alllanglex and alllangdir 
In 24 of the 60 Chadic languages which have a directional extension there is evidence of a directional extension being used in idiomatic expressions---the verb and extension form a predicate whose meaning is not composed from the sum of its parts.\footnote{The 24 languages include languages where the only attested idiomatic use is to distinguish `buy' and `sell' meanings (Chapter \ref{ch:buysell}). If these were excluded, there would still be 19 languages where directional extensions have an attested idiomatic function. These numbers exclude the use of directional extensions with a non-motion verb to express directed caused accompanied motion (CAM). In Chapter \ref{ch:CAM} it will be discussed whether these CAM expressions are also idiomatic expressions. Including these cases, there would be 40 of 60 languages where directional extensions have an idiomatic function.} An overview of how many languages show evidence of idiomatic combinations by type of directional extension is shown in Table \ref{tab:DIR_lex}. For every type of directional extension included in this study, there is at least one language where that same morpheme has an idiomatic function. Two particularly common idiomatic meanings are discussed in detail in following chapters: caused accompanied motion (Chapter \ref{ch:CAM}) and buying-selling events (Chapter \ref{ch:buysell}). Otherwise, there are not many clear semantic trends in the available data. 

%data in tab:DR_plustable2 is from Python code in Functions 
%Counted as OUT and VENT (Dghwede)
%VENT.OUT &    1 &             1 \\
\begin{table}
	\caption{Number of Chadic languages in which each type of directional extension is attested with an idiomatic use}
	\label{tab:DIR_lex}
	\begin{tabular}{lrr} % add l for every additional column or remove as necessary
		\lsptoprule
      Path &  DIR &  +Idiom \\
\midrule
VENT &   47 &             7 \\
ITIVE &   19 &             9 \\
INTO &   16 &             3 \\
OUT &   15 &             7 \\
UP &    8 &             5 \\
DOWN &   13 &             3 \\
ONTO &    9 &             6 \\
UNDER &    3 &             2 \\
HOME &    2 &             1 \\
BUSH &    1 &             1 \\
SIDE &    1 &             1 \\
		\lspbottomrule
	\end{tabular}
\end{table}

Although ventive directional extensions are the most common type of directional extension, they are not commonly found in idiomatic uses (not including CAM). Even within the seven languages that are counted as using a ventive directional extension idiomatically, the available examples are limited. In four languages (\hyperref[bole]{Bole}, \hyperref[gizi]{Giziga}, \hyperref[molo]{Moloko} and \hyperref[soko]{Sokoro}), the idiomatic use of ventive is only atttested in the context of making a distinction between buying and selling events (Chapter \ref{ch:buysell}). In the two Hurza (Biu-Mandara) languages, \hyperref[vame]{Vamé} and \hyperref[mbuk]{Mbuko}, the same extension that expresses ventive direction is used with a non-motion verb root glossed `be' to express ventive motion, `come', as in (\ref{ex:mbuk3}). 

\ea\label{ex:mbuk3}
\langinfo{\hyperref[mbuk]{Mbuko} verb \textit{n-} `be' with ventive suffix}{}{\cite[8]{gravinaVerbPhraseMbuko2001}} \\
\gll a-n-ay ahay a-juvok \\	
3\gl{sg}.\gl{pfv}-be-\gl{vent} to.here to-guest.hut	\\
\glt `He came to the guest hut.'
\z 

In the case of \hyperref[gude]{Gude} there is just one example of an idiomatic use of the ventive directional extension, shown in (\ref{ex:gude0b}). Notably, the verb in this example is a manner-of-motion verb. 

\ea\label{ex:gude0b}
\langinfo{\hyperref[gude]{Gude}  verb \textit{hwa} `run' with ventive suffix}{}{\cite[203]{hoskinsonGrammarDictionaryGude1983}} \\
\gll hwa-ya	\\ 
run-\gl{vent}	\\ 
\glt `grow horizontally (of vine)'
\z 

In \hyperref[dghw]{Dghwede}, one example of the outward-ventive extension, shown in (\ref{ex:dghw0}), could be interpreted as a motion-related metaphor or as an idiomatic collocation. 

\ea\label{ex:dghw0}
\langinfo{\hyperref[dghw]{Dghwede} verb \textit{pày} `show' with ventive outward suffix}{}{\cite[23]{frickVerbalSystemDghwede1978}} \\
\gll ngá pày-awíle páyá và	\\
?? show-\gl{vent}.\gl{out} show ??	\\
\glt `It did not show forth (was not obvious).'
\z 

Overall, given that only seven of 47 languages with a ventive directional extension have any evidence of idiomatic usage, it appears that non-ventive directional extensions are more prone to taking on fixed idiomatic meanings than ventive directional extensions. This is another way in which ventive directional extensions differ from other directional extensions in Chadic languages. 

In general, the meanings of attested idiomatic combinations vary considerably, as illustrated in (\ref{ex:glav0}), (\ref{ex:malg02}), (\ref{ex:bura02}) and (\ref{ex:psik24}). At least from an outsider's perspective, the meanings of these combinations of verb and verbal extension are not predictable from their parts, and the event described is often a completely different type of event than that normally predicated by the verb. 

\ea\label{ex:glav0}
\langinfo{\hyperref[glav]{Glavda} verb \textit{xən} `spend the day' with upward suffix}{}{\cite[]{owensGlavdaRuralUrban2011}} \\
\gll xən-it	\\
spend.day-\gl{up}		\\
\glt `spend the night'

\ex\label{ex:malg02}
\langinfo{\hyperref[malg]{Malgwa} verb \textit{ja} `hit' with upward suffix}{}{\cite[166]{lohrSpracheMalgwaNara2002}} \\
\gll ja-n-an-t-áá-və́-je	\\
hit-1\gl{sg}.\gl{prf}-\gl{obj}-\gl{tr}-\gl{up}-\gl{appl}-hit	\\
\glt `I walked around.' [German: `\textit{Ich bin herumgegangen}.']

\ex\label{ex:bura02}
\langinfo{\hyperref[bura]{Bura} verb \textit{dzib} `fill with food' with `side' suffix}{}{\cite[4]{blenchBuraVerbalExtensions2010}} \\
\gll dzib-dza	\\
fill.w.food-\textsc{side} \\
\glt `to plaster a mat house'

\ex\label{ex:psik24}
\langinfo{\hyperref[psik]{Psikye} verb \textit{mene} `do' with outward suffix}{}{\cite[114]{smithKapsikiLanguage1969}} \\
\gll  'a ké-mene-vé tlené	\\
\textsc{active} \gl{ptcp}-do-\gl{out} work	\\
\glt `He did the work for wages.'
\z
 
It is also reported in several languages that there are certain verb roots that do not occur without a verbal extension. For example, in \hyperref[shaa]{Ron-Scha}, there are three verb roots for which the ventive suffix \textit{-o} is obligatory, as in (\ref{ex:shaa4}). The verb roots in these cases are not known to have any other function. These appear to be lexicalized combinations, but the verb root without the verbal extension has gone out of use. 
 
 \ea\label{ex:shaa4}
 \langinfo{\hyperref[shaa]{Ron-Scha} verb \textit{wo'} `wake up' with (obligatory) ventive suffix}{}{\cite[269]{jungraithmayrRonSprachenTaschadohamitischeStudien1970}} \\
 \gll wo'-ó(h) \\
 ??-\gl{vent} \\
 \glt `wake up' [German: `{``}\textit{her}''-\textit{aufwecken}']
 \z 

\section{Conclusion}

It is not surprising that verbal extensions expressing direction also have non-directional functions when used with non-motion verbs. A similar pattern has been discussed extensively in regard to verbal particles with a directional meaning in English and other European languages in constructions known as phrasal verbs or verb-particle constructions \citep{deheVerbparticleExplorations2002,luoParticleVerbsEnglish2019}. In terms of morphosyntax, directional extensions in Chadic languages appear to be less complex than English verb-particle constructions, but semantically there are many similarities. For example, the particle \textit{up} in English can have a directional meaning, as in \textit{toss up (in the air)}, an aspectual meaning, as in \textit{eat up}, or appear in idiomatic verb-particle combinations, as in \textit{blow up} (= \textit{explode}) \citep{jackendoffEnglishParticleConstructions2002}.\footnote{Outside of Chadic and Germanic languages, idiomatic combinations of a verb root and directional verbal morphology are also attested in Nilotic languages \citep[728]{payneExtensionAssociatedMotion2021}.}

While the similarities are apparent, there is also a major difference. Directional verbal particles in English do not co-express associated motion. Among Chadic languages, associated motion, especially subsequent motion, is the most common meaning to find co-expressed in a directional extension. Co-expression of associated motion with (ventive) directional meaning is also commonly found in other families of African languages, including in Berber languages and Akkadian (Semitic) within the Afroasiatic phylum \citep{belkadiAssociatedMotionDeictic2015,kovalAssociatedMotionOld2024}.

From a global perspective, there is no clear preference for or against the co-expression of associated motion and direction. In a cross-linguistically balanced sample of 114 languages with directional morphemes, half of the directional morphemes are found to co-express associated motion\index{associated motion} \citep{rossCrosslinguisticSurveyAssociated2021,rossPseudocoordinationSerialVerb2021}. Globally, the co-expression of associated motion appears to be more common with ventive directional extensions. For example, in languages with vertical directional extensions, \citet[281]{rossPseudocoordinationSerialVerb2021} finds that only six of 50 (12\%) co-express associated motion. Therefore, it may be the lack of ventive directional particles in English that helps explain the absence of the co-expression of associated motion. 

There is also variation within Chadic languages regarding what meanings tend to be co-expressed with which types of directional extensions. Nevertheless, some patterns do emerge. In particular, the additive interpretation (Section \ref{sec:additive}) appears to be found co-expressed only with `onto' directional extensions, and the `split' interpretation (Section \ref{sec:objectsplit}) appears only with inward directional extensions. Beneficiary interpretations (Section \ref{sec:beneficiary}) are most common with ventive directional extensions, but are also attested with the \hyperref[malg]{Malgwa} outward directional extension. While ventive directional extensions are commonly found to co-express a beneficial meaning and/or subsequent associated motion, they are less likely to be found in an idiomatic verb-extension combination compared to other types of directional extensions. Another area of research to explore regarding the multifunctionality of directional extensions is the role that the lexical semantics of the verb root plays in influencing how a directional extension will be interpreted. This question is explored in the next chapter. 

%Berber languages also exhibit aspectual or adverbial meaning co-expressed with directional meaning in verbal morphology \citep{mettouchiGrammaticalizationDirectionalClitics2011,belkadiDeicticDirectionalitySpace2015}. For example, \citet[55]{kossmannGrammaticalSketchGhadames2013a} describes uses of directional morphology in Ghadames Berber including marking the relative distance between speaker and hearer with verbs of saying or marking events that include a change from invisibility to visibility. 

%The pattern of co-expression of direction and aspectual meaning is also found in other languages around the world. For example, \citet{hooperDeixisAspectTokelauan2002} presents a comprehensive overview of aspectual uses of directional particles in Tokelauan. A similar pattern can be found in Australian Kriol where an upward directional suffix can also express that ``the action carried out to its fullest extent'' and a downward directional suffix can also express ``cessation of one state and at the same time, change to a different state'' \citep[43--44]{hudsonGrammaticalSemanticAspects1981}. 
